%% Commands for TeXCount
%TC:macro \cite [option:text,text]
%TC:macro \citep [option:text,text]
%TC:macro \citet [option:text,text]
%TC:envir table 0 1
%TC:envir table* 0 1
%TC:envir tabular [ignore] word
%TC:envir displaymath 0 word
%TC:envir math 0 word
%TC:envir comment 0 0
%%
%%
%% The first command in your LaTeX source must be the \documentclass command.
\documentclass[acmsmall,review,screen]{acmart}
\settopmatter{}
%%
%% \BibTeX command to typeset BibTeX logo in the docs
%\AtBeginDocument{%
%  \providecommand\BibTeX{{%
%    Bib\TeX}}}
\startPage{1}

%\usepackage{txfonts}
\usepackage{subcaption} %% For complex figures with subfigures/subcaptions
%% http://ctan.org/pkg/subcaption
\usepackage{uri}
\usepackage{mathtools}
\usepackage{amsmath}
\usepackage{amsthm}

\clubpenalty = 10000
\widowpenalty = 10000
\displaywidowpenalty = 10000


\hyphenation{Comp-Cert}

\newcommand{\TODO}[1]{\textbf{\textcolor{red}{[ TODO: #1]}}}

%%
%% Submission ID.
%% Use this when submitting an article to a sponsored event. You'll
%% receive a unique submission ID from the organizers
%% of the event, and this ID should be used as the parameter to this command.
%%\acmSubmissionID{123-A56-BU3}

%%
%% For managing citations, it is recommended to use bibliography
%% files in BibTeX format.
%%
%% You can then either use BibTeX with the ACM-Reference-Format style,
%% or BibLaTeX with the acmnumeric or acmauthoryear sytles, that include
%% support for advanced citation of software artefact from the
%% biblatex-software package, also separately available on CTAN.
%%
%% Look at the sample-*-biblatex.tex files for templates showcasing
%% the biblatex styles.
%%

%%
%% The majority of ACM publications use numbered citations and
%% references.  The command \citestyle{authoryear} switches to the
%% "author year" style.
%%
%% If you are preparing content for an event
%% sponsored by ACM SIGGRAPH, you must use the "author year" style of
%% citations and references.
%% Uncommenting
%% the next command will enable that style.
%%\citestyle{acmauthoryear}



%%
%% end of the preamble, start of the body of the document source.
\begin{document}

%%
%% The "title" command has an optional parameter,
%% allowing the author to define a "short title" to be used in page headers.
\title{A Formal Model of Spot Markets in Rocq}

\author{William Mansky}
\orcid{0000-0002-5351-895X}
\affiliation{%
	\institution{University of Illinois Chicago}
	\city{}
	\country{USA}
}
\email{mansky1@uic.edu}

%etc.

%%
%% By default, the full list of authors will be used in the page
%% headers. Often, this list is too long, and will overlap
%% other information printed in the page headers. This command allows
%% the author to define a more concise list
%% of authors' names for this purpose.
%\renewcommand{\shortauthors}{Trovato et al.}

%%
%% The abstract is a short summary of the work to be presented in the
%% article.
\begin{abstract}
\end{abstract}

%%
%% The code below is generated by the tool at http://dl.acm.org/ccs.cfm.
%% Please copy and paste the code instead of the example below.
%%

\begin{CCSXML}
	<ccs2012>
	<concept>
	<concept_id>10011007.10011074.10011099.10011692</concept_id>
	<concept_desc>Software and its engineering~Formal software verification</concept_desc>
	<concept_significance>500</concept_significance>
	</concept>
	</ccs2012>
\end{CCSXML}

\ccsdesc[500]{Software and its engineering~Formal software verification}

%%
%% Keywords. The author(s) should pick words that accurately describe
%% the work being presented. Separate the keywords with commas.
\keywords{}

%%
%% This command processes the author and affiliation and title
%% information and builds the first part of the formatted document.
\maketitle

\section{An Abstract Spot Market}
A spot market is defined over a set of participants $\mathcal{P}$ and a set of resources $\mathcal{R}$. The state of the market is a pair $(A, p)$ where $A : \mathcal{R} \rightharpoonup \mathcal{P}$ is an allocation of resources to participants, and $p : \mathbb{N}$ is the current price of the resources. (For this simple model, we assume that all resources are interchangeable and have the same price.) We also assume a fixed map $B : \mathcal{P} \rightarrow \mathbb{N}$ from each participant to their \emph{bid}, the amount they are willing to pay for a resource. (In later iterations we will allow $B$ to change over time.)

The market evolves through a sequence of \emph{steps} $(A, p) \rightarrow (A', p')$. Conceptually, in each step:
\begin{enumerate}
\item The seller chooses a new price for the resources.
\item Participants who are not willing to pay the new price relinquish their resources.
\item Available resources are allocated to participants who are willing to pay for them.
\end{enumerate}
In this model we assume that the combination of these three operations is \emph{atomic}, i.e., there is no observable state of the market in which some but not all of these operations have occurred. We will refine this assumption later.

Formally:
\begin{align*}
&\mathsf{drop}(p, A) \triangleq \{(r, a).\ (r, a) \in A \land p \leq B(a)\}\\
&\mathsf{interested}(p, A, a) \triangleq p \leq B(a) \land \forall r.\ A(r) \neq a\\
&\mathsf{buy}(p, A, A') \triangleq \exists C.\ A' = A \uplus C \land (\forall r, a.\ C(r) = a \Rightarrow \mathsf{interested}(p, A, a))\ \land \\&\qquad (\forall r.\ r \not\in \mathsf{dom}(A') \Rightarrow \forall a.\ \neg \mathsf{interested}(p, A, a))\\
&(A, p) \rightarrow (A', p') \triangleq \mathsf{buy}(p', \mathsf{drop}(p', A), A')
\end{align*}
The new price is chosen arbitrarily, reflected in the fact that $\rightarrow$ imposes no restrictions on $p'$; if we want to analyze a more concrete market, we can fill in a pricing strategy. The $\mathsf{drop}$ operation removes all entries from $A$ whose participants' bids are lower than the new price. The $\mathsf{buy}$ operation arbitrarily allocates unused resources by choosing a new map $C$ such that 1) all newly assigned resources are assigned to interested participants, and 2) a resource is only unused if no interested participants remain, where an \emph{interested} participant is one who is willing to pay the (new) price and not already assigned a resource.
Although the $\mathsf{buy}$ operation is nondeterministic, there is always at least one way to assign resources to interested parties:
\begin{theorem} For any $p$ and $A$, there exists an $A'$ such that $\mathsf{buy}(p, A, A')$.
\end{theorem}

One desirable property of spot markets is \emph{monotonicity}: increasing the price should not increase the resources allocated, and decreasing the price should not decrease the resources allocated. The model of this section is guaranteed to satisfy these properties.
\begin{theorem} If $(A, p) \rightarrow (A', p')$ and $p \leq p'$, $\mathsf{dom}(A') \subseteq \mathsf{dom}(A)$.
\end{theorem}
\begin{theorem} If $(A, p) \rightarrow (A', p')$ and $p' \leq p$, $A \subseteq A'$.
\end{theorem}

\section{A Less Atomic Spot Market}
The assumptions of the previous model rule out many attacks that are possible in real markets. As one example, the assumption of atomicity rules out a denial-of-service attack where brief price spikes force participants to leave and re-enter the market, so that most of their resource-time is spent in shutdown and startup rather than utilization. Our second model exposes this attack vector by introducing a delay between the announcement of a new price and its implementation (analogous to Amazon's Spot Instance interruption notices), and a startup delay on newly allocated resources. We also allow the bids $B$ to change over time, so that an attacker may briefly submit a very high bid (to try to cause a price spike) and then withdraw it.

In this model, each allocated resource is in a state $s \in \mathcal{S} \triangleq \mathsf{Starting}(a) \ |\ \mathsf{Used}(a)$, where $a \in \mathcal{P}$ is its assigned participant; in either case we write $\mathsf{user}(s) = a$. The state of the market is a tuple $(A, p, q, B)$ where $A : \mathcal{R} \rightharpoonup \mathcal{S}$ is the allocation, $p : \mathbb{N}$ is the current price, $q : \mathbb{N}^?$ is an optional incoming price, and $B : \mathcal{P} \rightarrow \mathbb{N}$ is the current bids of the participants. A market step $(A, p, q, B) \rightarrow (A', p', q', B')$ is now a combination of four operations:
\begin{enumerate}
\item The seller either announces a new price, or implements the previously announced price.
\item Participants who are not willing to pay the announced price relinquish their resources.
\item Previously allocated resources start up, moving from $\mathsf{Starting}$ to $\mathsf{Used}(a)$.
\item Available resources are allocated to participants who are willing to pay for them.
\end{enumerate}

Formally:
\begin{align*}
&\mathsf{set\_price}(p, q, p', q') \triangleq (q = p' \land q' = \bot) \lor (q = \bot \land p' = p)\\
&\mathsf{drop}(p, B, A) \triangleq \{(r, s).\ (r, s) \in A \land p \leq B(\mathsf{user}(a))\}\\
&\mathsf{start}(A) \triangleq \{(r, s').\ (r, s) \in A \land \text{if } s = \mathsf{Starting}(a) \text{ then } s' = \mathsf{Used(a)} \text{ else } s' = s\}\\
&\mathsf{interested}(p, B, A, a) \triangleq p \leq B(a) \land \forall r.\ \mathsf{user}(A(r)) \neq a\\
&\mathsf{buy}(p, B, A, A') \triangleq \exists C.\ A' = A \uplus \{(r, \mathsf{Starting}(a).\ (r, a) \in C\}\ \land \\&\qquad (\forall r, a.\ C(r) = a \Rightarrow \mathsf{interested}(p, B, A, a)) \land (\forall r.\ r \not\in \mathsf{dom}(A') \Rightarrow \forall a.\ \neg \mathsf{interested}(p, B, A, a))\\
&(A, p, q, B) \rightarrow (A', p', q', B') \triangleq \mathsf{set\_price}(p, q, p', q') \land \mathsf{buy}(p', B', \mathsf{start}(\mathsf{drop}(q', B', A)), A')
\end{align*}
If a price $q$ has been announced, $p'$ is constrained to be $q$ and no new price is announced; otherwise ($q = \bot$) the price is unchanged but a new price may be announced arbitrarily. Previously bought resources transition from $\mathsf{Starting}$ to $\mathsf{Used}$, and newly bought resources begin in the $\mathsf{Starting}$ state. Participants drop based on the newly announced price (if any), not the current price.

\end{document}
